\section{Tessuto connettivo propriamente detto}

I tipi di tessuti connettivi che costituiscono questa categoria differiscono tra loro in base alla proporzione, variabile da un tipo all'altro, di \textbf{cellule} -- \textbf{fibre} -- \textbf{sostanza fondamentale}.

\begin{itemize}
  \item tessuto connettivo fibroso lasso;
  \item tessuto connettivo fibroso denso;
  \item tessuto elastico;
  \item tessuto reticolare;
  \item tessuto adiposo.
\end{itemize}

% ---------------------------------------------------------------------------
\subsection{Tessuto connettivo fibroso lasso}

Costituisce il tessuto ``d'imballaggio'' del corpo.
\begin{itemize}
  \item prevale la sostanza fondamentale;
  \item scarse cellule (fibroblasti, macrofagi, adipociti);
  \item scarse fibre collagene;
  \item frequenti fibre elastiche;
  \item vascolarizzazione marcata;
  \item innervazione presente.
\end{itemize}

\rubrica{Funzioni}
\begin{itemize}
  \item fornisce protezione e ammortizza traumi per organi ed apparati;
  \item supporto per l'epitelio del tratto gastroenterico, respiratorio e urinario;
  \item garantisce ancoraggio e vascolarizzazione per il tessuto epiteliale;
  \item funzione di sostegno e protezione attorno ad alcuni organi (es.\ rene);
  \item funzione trofica, in quanto, penetrando tra gli organi, costituisce l'ambiente in cui decorrono i vasi sanguigni;
  \item interviene nei processi di riparazione dei danni tissutali producendo fibre che portano alla formazione di un tessuto cicatriziale.
\end{itemize}

% ---------------------------------------------------------------------------
\subsection{Tessuto connettivo fibroso denso}

Detto anche \textbf{tessuto collageno}, in quanto costituito in prevalenza da fibre collagene. Le fibre possono essere orientate parallele le une alle altre (tessuto denso regolare) o ad intreccio (tessuto denso irregolare).
\begin{itemize}
  \item cellule presenti: fibrociti;
  \item sostanza fondamentale scarsissima;
  \item vascolarizzazione scarsa;
  \item innervazione presente;
  \item altamente resistente alle sollecitazioni meccaniche.
\end{itemize}
Esempio: tendini e fasce muscolari, capsule di organi parenchimatosi.

\rubrica{Tessuto denso regolare (o a fasci paralleli)}
Fibre collagene strettamente impacchettate che decorrono parallelamente tra loro.
\begin{itemize}
  \item scarsa componente amorfa;
  \item scarsa componente cellulare, costituita da fibrociti allungati e disposti in file regolari tra i fascetti di fibre collagene;
  \item notevole resistenza alla trazione;
  \item tipico dei tendini.
\end{itemize}

\rubrica{Tessuto denso irregolare (o a fasci incrociati)}
Costituito da lamelle sovrapposte in cui le fibre sono tutte parallele e incrociano con angoli diversi le fibre di lamelle contigue.
\begin{itemize}
  \item i fibrociti hanno forma allungata e appiattita;
  \item esempio: derma.
\end{itemize}

\rubrica{Funzioni}
\begin{itemize}
  \item fornitura di solido ancoraggio;
  \item trasmissione della contrazione muscolare;
  \item stabilizzazione delle relative posizioni delle ossa;
  \item riduzione dell'attrito tra muscoli;
  \item capacità di resistenza a forze applicate da molte direzioni.
\end{itemize}

% ---------------------------------------------------------------------------
\subsection{Tessuto elastico}

È un tessuto connettivo denso regolare in cui le fibre sono quasi esclusivamente fibre elastiche $\rightarrow$ fortemente estensibili, fino al 150\% della lunghezza iniziale. Al cessare della forza che ha indotto l'estensione, sono in grado di riprendere le dimensioni originarie.

\rubrica{Funzioni}
\begin{itemize}
  \item garantisce la stabilizzazione della posizione;
  \item consente l'ammortizzamento dei traumi;
  \item rende possibile l'espansione e la contrazione di organi.
\end{itemize}

Il tessuto connettivo elastico forma la tonaca media dei grossi vasi arteriosi, quali l'aorta.

% ---------------------------------------------------------------------------
\subsection{Tessuto reticolare}

Costituito in prevalenza da fibre reticolari.
\begin{itemize}
  \item cellule scarse (fibroblasti, macrofagi);
  \item vascolarizzazione e innervazione scarsa;
  \item sostegno organizzativo di organi ghiandolari (fegato, rene) e non (milza, linfonodi, tessuto emopoietico).
\end{itemize}

% ---------------------------------------------------------------------------
\subsection{Tessuto adiposo}

Costituito in prevalenza da cellule (adipociti) cariche di lipidi.
\begin{itemize}
  \item sostanza fondamentale scarsa;
  \item fibre scarse (reticolari);
  \item vascolarizzazione moderata;
  \item innervazione moderata;
  \item metabolismo cellulare marcato, con continua scissione e sintesi di lipidi.
\end{itemize}

\rubrica{Funzioni}
\begin{itemize}
  \item riserva di materiali energetici $\rightarrow$ funzione trofica;
  \item sistema di rivestimento coibente $\rightarrow$ evita la dispersione del calore interno;
  \item protezione meccanica e di sostegno;
  \item 50\%: \textbf{tessuto adiposo di copertura}, localizzato nel pannicolo sottocutaneo, con funzione coibente e meccanica;
  \item 45\%: \textbf{tessuto adiposo interno}, dislocato nella cavità addominale;
  \item 5\%: \textbf{grasso di infiltrazione}, localizzato nel tessuto muscolare, dove agevola la funzione biomeccanica dei muscoli.
\end{itemize}

\rubrica{Classificazione funzionale}
\begin{itemize}
  \item \textbf{tessuto di deposito}: varia in relazione allo stato di nutrizione dell'organismo;
  \item \textbf{tessuto di sostegno}: non è soggetto a variazioni quantitative e si trova nella pianta del piede e della mano;
  \item \textbf{tessuto adiposo uniloculare}: presente nell'uomo adulto e nella maggior parte dei mammiferi;
  \item \textbf{tessuto adiposo multiloculare}: frequente nei bambini, nei piccoli mammiferi e negli animali ibernanti.
\end{itemize}

\rubrica{Tessuto adiposo uniloculare (o univacuolare, o tessuto adiposo bianco o giallo)}
\begin{itemize}
  \item costituito da cellule a stretto contatto, con scarsa matrice extracellulare;
  \item adipociti molto voluminosi (diametro fino a 150\,$\mu$m);
  \item presenza di una grossa gocciola lipidica;
  \item citoplasma molto ridotto;
  \item nucleo confinato alla periferia della cellula, compresso contro la membrana plasmatica;
  \item presenza, intorno alla cellula, di un involucro glicoproteico e fibre reticolari disperse in una scarsa componente amorfa.
\end{itemize}

\rubrica{Trasporto dei lipidi tra capillare e adipocita}
I lipidi sono trasportati nel torrente circolatorio sotto forma di chilomicroni e di lipoproteine a bassissima densità (VLDL). L'enzima lipoprotein lipasi, elaborato dall'adipocita e penetrato nel lume del capillare, scinde i lipidi in acidi grassi e glicerolo. Gli acidi grassi si diffondono nel connettivo del tessuto adiposo ed entrano negli adipociti, dove vengono di nuovo esterificati in trigliceridi di accumulo. Quando necessario, i trigliceridi accumulati vengono scissi da una lipasi ormone-dipendente in acidi grassi e glicerolo; queste sostanze passano nel connettivo del tessuto adiposo e, da qui, nei capillari, dove verranno legate all'albumina e trasportate col sangue. Il glucosio può essere trasportato dal sangue agli adipociti che lo possono trasformare in grasso di riserva.

\rubrica{Tessuto adiposo multiloculare (o multivacuolare, o tessuto adiposo bruno)}
\begin{itemize}
  \item cellule notevolmente più piccole;
  \item lipidi presenti in numerose microgocce disperse in tutto il citoplasma;
  \item presenza di numerosi mitocondri;
  \item sulla membrana mitocondriale interna è presente la \textbf{termogenina}, particolare proteina transmembrana che funziona come canale per i protoni che si muovono dallo spazio intermembrana alla matrice mitocondriale, generando calore che viene ceduto al sangue circolante presente nella ricca rete vascolare del tessuto;
  \item tipico degli animali ibernanti.
\end{itemize}

\rubrica{Fisiologia del tessuto adiposo}
Metabolismo molto attivo, con un continuo cambio e scambio tra tessuti di riserva e bisogno energetico. Coinvolto nei processi metabolici come appetito, fertilità, difesa immunitaria. Mantiene costante il numero di cellule, ma può variare il contenuto di lipidi di ogni singola cellula.
